% Pursuing Stacks (cotes 134-2 to 134-8): archive cote 134-7, site page 38 (PDF page 39), typescript page 480.
% typescript. Draft compared: ps6.tex:1836-1904.
% AI-transcribed by Claude Opus 5.5 (claude-opus-5-5, adaptive thinking, effort high), Message Batches API, batch ps3 (msgbatch_01XdFmuVVBAcgc7EFxdtCqsj), 2026-11-12; not yet checked by eye.
\scanpage{134-7/38}
\mspage{480}

At this point it should be noted that the devissage we'll have to use
here is the "brutal" one,  which we frowned upon earlier today !  To
develop such a formalism,  it seems essential to work with smooth Kan
complexes and projective resolutions of the\inserted[typed]{k-modules} $\hand{\pi}_{\unsure{i|1}}$ as they appear one by one,
\xout{endowed with}  Whether we want to describe\xout{soft} "hard" or "soft" Postnikov
dévissage (see p.468 for the latter),  one common key step is the
\uline{linearization map}  coming from the exponential map (K)  applied component-%
wise
\[
(W) \qquad X_{\ast} \longrightarrow \xout{L(X_{\ast})}\ W^{\wedge}(L(X_{\ast})) \ ,
\]
which we would like to look upon as \xout{being} defining a homotopy class of
maps in $\hand{\mathcal{M}}_n(k)$
\[
\margin{X}\ (W\ ?) \qquad X_{\ast} \longrightarrow W(L_{pt}(X_{\ast})) \quad ,
\]
where the second member moreover is endowed with it's natural abelian
group structure (it's components are abelian group objects of U(k) and the
simplicial maps are additive).  To pass from (W) to (W ?) ,  it is felt
that the essential step is that (V) above be a quasi-isomorphism,\xout{x} hence\hand{,}
applying the functor W\hand{,} we should get a \xout{quasi-isomorphism} weak equivalence\hand{,}
hence an isomorphism in the \xout{suitable} derived category Hot$_n$(k),  hence
(W)  implies (W ?).  The \xout{main} main \deleted{trouble here}\inserted{flaw in this ``argument''} comes from the W$^{\wedge}$ in the
second member of (W),  which is'nt quite the same as W definitely.  Thus
some\inserted[typed]{further} amount of work will be needed presumably to get (W ?) from (W).
Of course,  we cant possibly just keep (W) as it is,  as for getting
devissage we need a map in $\hand{\mathcal{M}}_n$(k),  whereas the map (W) is\inserted{just} a map of
semisimplicial sheaves on Aff$_{/k}$  , \deleted{but}\inserted{where} the second member is \uline{not} in
$\hand{\mathcal{M}}_n$(k)  i.e. it's components are not in U(k).  Once we \unsure{got} (W ?)
factoring (W) up to homotopy (NB of course we assume $X_{\ast}$ to be an sK
complex in all this),  we still need a reasonable notion of homotopy
fibers of maps in $\hand{\mathcal{M}}_n$(k),  in order to push through the inductive step.

\qquad Thus,  a large part of the weight of the work ahead may well lie
upon \uline{developping the standard homotopy constructions} \uline{within the model}
\uline{category} $\hand{\mathcal{M}}_n$(k),  as contemplated on page 464.  This should be fun,  if
it can be done indeed\over{.}{!}  One difficulty here seems to be that \xout{there} Quil-%
len's  standard machines  wo'nt work,\inserted{not} "telles quelles" at any rate,  becau-%
se of the category U(k) failing to be stable under finite limits - it
does'nt even have fi\over{t}{b}er products.  But I think I'll stop my ponderings
on schematization here ...

\margin{38}
